\section{Introduction}
\label{sec:intro}

School spirit is a vital component of campus culture, fostering a sense of community and belonging among students, faculty, and alumni. At the University of North Carolina at Charlotte (UNCC), this spirit is often visibly manifested through the wearing of university-branded merchandise. However, quantifying this "school spirit" in real-time during campus events remains a challenge. Traditional methods of manual counting are tedious, error-prone, and unscalable. To address this, we propose the Charlotte Merchandise Model (CMM), an automated computer vision system capable of detecting and counting UNCC merchandise in live video streams.

The core of our solution leverages the power of deep learning for object detection. Specifically, we utilize the You Only Look Once (YOLO) architecture \cite{redmon2016you}, which has revolutionized real-time object detection by framing it as a single regression problem. By training a model to recognize specific UNCC logos and associated apparel (torsos and headgear), we aim to create a "School Spirit Meter" that can provide immediate feedback on crowd engagement.

Detecting specific brand logos in the wild presents unique challenges, including variations in lighting, occlusion, and the presence of visually similar but non-affiliated merchandise. As noted in recent surveys on logo detection \cite{hou2022deep}, while general object detectors are powerful, adapting them for specific logo recognition requires careful dataset curation and pipeline design. Our work builds upon these foundations to create a tailored solution for the UNCC community.

\section{Related Work}
\label{sec:related}

The foundation of our work lies in the advancements of real-time object detection, particularly the YOLO family of models. Redmon et al. introduced YOLO \cite{redmon2016you} as a unified, real-time object detection system. Unlike prior region-based approaches, YOLO treats detection as a single regression problem, predicting bounding boxes and class probabilities directly from full images in one evaluation. The original YOLO achieved 63.4\% mAP on the PASCAL VOC dataset at 45 FPS, making it the fastest detector of its time. However, it historically struggled with small or clustered objects due to spatial grid constraints. This foundational work established the viability of single-shot detection, which is critical for our real-time application.

In the domain of logo detection, Hou et al. \cite{hou2022deep} provide a comprehensive survey, highlighting that YOLO-based architectures are strong candidates for logo detection tasks. They discuss various pipelines, including YOLO-to-classifier approaches, which can offer higher accuracy at the cost of speed, versus single-shot detectors like YOLO which are optimized for "on-the-fly" classification. The survey notes that while two-stage detectors might achieve higher precision, the latency introduced is often unacceptable for live video processing. Given our requirement for a live "School Spirit Meter," we prioritized the single-shot approach for its real-time capabilities, accepting a potential trade-off in absolute precision for the sake of interactivity.

Recent research continues to improve crowd detection, a relevant aspect of our project. Wen et al. \cite{wen2024research} explored enhancing YOLOv5 with Focus layers and Convolutional Block Attention Modules (CBAM) to better handle dense crowds and small objects. Their work suggests that attention mechanisms can significantly improve accuracy in complex scenes, which informs our understanding of the challenges in detecting merchandise in crowded campus events. Specifically, the Focus layer reduces computation while increasing speed, which is beneficial for live detection, while the CBAM module helps in focusing on small windows and spots of specified merchandise, improving accuracy in dense crowds. These insights are particularly relevant as we aim to deploy our model in bustling campus environments.


